\documentclass[11pt]{article}

\usepackage[margin=1in]{geometry}
\usepackage[T1]{fontenc}
\usepackage[utf8]{inputenc}
\usepackage{amsmath,amssymb}
\usepackage{float}
\usepackage{graphicx}
\usepackage{xcolor}
\usepackage[
    colorlinks=true,
    linkcolor=blue!60!black,
    citecolor=blue!60!black,
    urlcolor=blue!60!black
]{hyperref}

\newcommand{\loss}{\mathcal{L}}
\newcommand{\R}{\mathbb{R}}
\newcommand{\img}[2][]{\includegraphics[#1]{\detokenize{#2}}}
\newcommand{\panelcaption}[1]{\par\vspace{0.35em}{\small #1}}

\title{Empirical Study of Large-Learning-Rate GD, Normalized GD, and Muon}
\author{}
\date{}

\begin{document}
\maketitle

\section{Experimental Setup}

\paragraph{Goal of the experiments.}
The experiments are designed to test whether the apparent advantage of Muon
could be explained away by an insufficiently tuned GD baseline.  The first
section gives GD favorable learning-rate schedules: constant steps scaled by
the initial Hessian stability threshold, polynomial decay, and line search.
The step-size trajectory of line-search GD is then used to motivate an
explicit warm-up-stable GD schedule.  The second section removes the Hessian
scaling altogether and sweeps very large absolute step sizes for GD and
normalized GD, so that the comparison is not tied to a particular stability
normalization.

\paragraph{Problem instance and objective.}
Each run uses $N=20000$ examples, dimensions $d\in\{64,128\}$, Zipf exponent
$\alpha=1.5$, initialization $W_0=0$, and $120$ optimization steps.  For
rank $i$,
\[
    u_i,v_i \overset{\mathrm{i.i.d.}}{\sim} \mathcal{N}(0,I_d/d),
    \qquad
    p_i=\frac{i^{-\alpha}}{\sum_{k=1}^N k^{-\alpha}}.
\]
For $W\in\R^{d\times d}$, define the score
$s_{ji}(W)=u_j^\top Wv_i$.  The loss is the full-softmax weighted
cross-entropy
\[
    \loss(W)
    =
    \sum_{i=1}^N p_i
    \left[
        \log\left(\sum_{j=1}^N \exp(s_{ji}(W))\right)
        -
        s_{ii}(W)
    \right].
\]
Thus every query is compared against all $N$ classes.  The plotted losses are
exact full-batch losses; no minibatching or negative sampling is used.

\paragraph{Optimizers.}
Let $g_t=-\nabla\loss(W_t)$.  All methods update
\[
    W_{t+1}=W_t+\eta_tD_t,
\]
but use different directions:
\[
    D_t^{\mathrm{GD}}=g_t,
    \qquad
    D_t^{\mathrm{NGD}}=\frac{g_t}{\|g_t\|_F},
    \qquad
    D_t^{\mathrm{Muon}}=\operatorname{Polar}(g_t).
\]
Here $\operatorname{Polar}(g_t)=UV^\top$ when
$g_t=U\Sigma V^\top$ is the compact SVD.  In other words, GD uses the raw
gradient magnitude, normalized GD keeps only the Frobenius-normalized gradient
direction, and Muon replaces the gradient by its spectral polar factor.

\paragraph{Learning-rate schedules in the first section.}
The first section asks whether a strong schedule for raw GD can match Muon.
To avoid handicapping GD with only small stable steps, define
\[
    \eta_{\mathrm{stab},0}
    =
    \frac{2}{\lambda_{\max}(H_0)},
    \qquad
    H_0=\nabla^2\loss(W_0).
\]
Constant GD uses $\eta_t=\rho\,\eta_{\mathrm{stab},0}$ with
\[
    \rho\in\{0.25,0.5,1,2,4,8,16\}.
\]
Decayed GD starts from a large Hessian-scaled step and then decreases it:
\[
    \eta_t
    =
    \frac{\rho\,\eta_{\mathrm{stab},0}}{(1+t/T)^p},
    \qquad
    \rho\in\{2,4,8,16\},\quad
    T\in\{20,80\},\quad
    p\in\{0.5,1,2\}.
\]
Exact line-search GD chooses $\eta_t$ along $D_t=g_t$ by approximately
minimizing $\loss(W_t+\eta D_t)$ over $\eta\ge0$, using geometric bracketing
with growth factor $2$ followed by $8$ bisection refinements.  The maximum
GD line-search step in this comparison is $64\,\eta_{\mathrm{stab},0}$.
The Muon baselines are constant-step Muon with
\[
    \eta\in\{0.25,0.5,1,1.5,2,3,4,5,6\}
\]
and line-search Muon along $D_t=\operatorname{Polar}(g_t)$ with maximum step
$32$.

\paragraph{Why add warm-up-stable GD.}
The line-search GD step-size plots are not only a baseline; they also reveal a
possible hand-designed schedule for GD.  Empirically, the selected line-search
steps rise quickly from the initialization scale and then stay large for much
of the trajectory.  To test whether this behavior, rather than the online line
search itself, is enough to improve GD, we add the parametric schedule
\[
    \eta_t
    =
    \eta_{\mathrm{start}}
    \left(
        \frac{\eta_{\mathrm{hold}}}{\eta_{\mathrm{start}}}
    \right)^{\min(t/W,1)^q},
    \qquad
    \eta_{\mathrm{start}}=\rho_{\mathrm{start}}\eta_{\mathrm{stab},0},
    \qquad
    \eta_{\mathrm{hold}}=\rho_{\mathrm{hold}}\eta_{\mathrm{stab},0}.
\]
Here the sweep uses
\[
    \rho_{\mathrm{start}}=1,\qquad
    \rho_{\mathrm{hold}}\in\{8,16,32,64,96,128\},\qquad
    W\in\{3,5,10\},\qquad
    q\in\{1,2\}.
\]
This schedule is intentionally favorable to GD: it imitates the warm-up and
large-hold pattern suggested by exact line search, but it does not use Muon or
any spectral update direction.

\paragraph{Large absolute-step sweep in the second section.}
The second section addresses a different objection: perhaps scaling GD by
$2/\lambda_{\max}(H_0)$ is itself too restrictive.  We therefore sweep
absolute learning rates directly.  Constant GD uses
\[
\begin{aligned}
    \eta\in\{&
    0.25,0.5,0.75,1,1.5,2,2.5,3,4,5,6,8,12,16,24,\\
    &32,48,64,96,128,192,256,384,512,768,1024,1536,\\
    &2048,3072,4096,6144,8192,12288\}.
\end{aligned}
\]
Normalized GD uses
\[
    \eta\in\{1,2,4,8,12,16,24,32,48,64,96\}.
\]
The same constant-step Muon grid is used as the baseline.  The main
large-step comparison also includes Muon line search and a small-step GD
line-search baseline capped at $\eta\le2.99$, which separates the small-step
regime from the aggressive large-step sweeps.

\paragraph{Selection rule and plotted quantities.}
The figures focus on two quantities: training loss and the selected step size
$\eta_t$.  Local sharpness diagnostics are also recorded along each
trajectory:
\[
    \widehat\lambda_{\max}(H(W_t)),
    \qquad
    \frac{2}{\widehat\lambda_{\max}(H(W_t))},
\]
estimated every $10$ steps by $3$ power iterations on a stratified subsample
of $2048$ queries.  For each method family and dimension, the displayed best
hyperparameter is selected by trajectory speed, not by final loss alone:
\[
    \operatorname{Score}(h)
    =
    \frac{1}{|\mathcal{S}|}
    \sum_{s\in\mathcal{S}}
    \frac{1}{121}
    \sum_{t=0}^{120}
    \log\bigl(\max\{\loss(W_t^{h,s}),10^{-30}\}\bigr).
\]
Unstable or missing trajectories receive a large penalty.  After this
selection step, plotted curves show the mean over available seeds with a
one-standard-deviation band.

\section{GD-Friendly Learning-Rate Schedules}

This section gives raw GD several increasingly favorable choices of step
size.  The first figure compares the best constant, decayed, and line-search
GD trajectories against constant and line-search Muon.  The next figure plots
the corresponding step sizes, which is used to diagnose what schedule the
line-search GD baseline is implicitly choosing.

\begin{figure}[H]
    \centering
    \img[width=0.9\linewidth]{figures/report_figures/schedule_loss_d64.png}
    \par
    \img[width=0.9\linewidth]{figures/report_figures/schedule_loss_d128.png}
    \caption{Loss comparison among constant Muon, exact-line Muon, constant GD, decayed GD, and exact-line GD. Top: $d=64$. Bottom: $d=128$.}
    \label{fig:schedule-loss}
\end{figure}

\begin{figure}[H]
    \centering
    \begin{minipage}{0.49\linewidth}
        \centering
        \img[width=\linewidth]{figures/report_figures/schedule_stepsize_d64.png}
        \panelcaption{(a) $d=64$. Selected or prescribed step size $\eta_t$ for each displayed method family.}
    \end{minipage}
    \hfill
    \begin{minipage}{0.49\linewidth}
        \centering
        \img[width=\linewidth]{figures/report_figures/schedule_stepsize_d128.png}
        \panelcaption{(b) $d=128$. Selected or prescribed step size $\eta_t$ for each displayed method family.}
    \end{minipage}
    \caption{Step-size trajectories for the selected schedules.}
    \label{fig:schedule-stepsize}
\end{figure}

The step-size curves motivate the warm-up-stable GD schedule in the setup:
line-search GD rapidly increases the step size from the initialization scale
and then keeps using a large step.  The following figure tests this
line-search-inspired schedule as an explicit GD baseline.

\begin{figure}[H]
    \centering
    \img[width=0.88\linewidth]{figures/report_figures/warmup_stable_loss_d64.png}
    \par\vspace{0.8em}
    \img[width=0.88\linewidth]{figures/report_figures/warmup_stable_loss_d128.png}
    \caption{Loss comparison after adding warm-up-stable GD. Top: $d=64$. Bottom: $d=128$.}
    \label{fig:warmup-loss}
\end{figure}

\section{Large Learning-Rate Sweeps and Normalized GD}

This section tests a more direct alternative: instead of parameterizing GD by
the initial Hessian scale, it sweeps absolute learning rates over a wide range.
The first two figures show the selected main comparisons.  The next two
figures expose the full constant-GD sweep against the selected Muon baseline,
and the last two figures repeat the same type of sweep for normalized GD.

\begin{figure}[H]
    \centering
    \img[width=0.88\linewidth]{figures/report_figures/large_lr_main_loss_d64.png}
    \caption{$d=64$. Main loss comparison including large constant-step GD, normalized GD, and Muon baselines.}
    \label{fig:large-sweep-loss-d64}
\end{figure}

\begin{figure}[H]
    \centering
    \img[width=0.88\linewidth]{figures/report_figures/large_lr_main_loss_d128.png}
    \caption{$d=128$. Main loss comparison under the same large-step sweep protocol.}
    \label{fig:large-sweep-loss-d128}
\end{figure}

\begin{figure}[H]
    \centering
    \img[width=0.88\linewidth]{figures/report_figures/gd_eta_sweep_vs_muon_best_d64.png}
    \caption{$d=64$. Constant-GD absolute learning-rate sweep compared with the selected constant Muon curve.}
    \label{fig:gd-eta-sweep-d64}
\end{figure}

\begin{figure}[H]
    \centering
    \img[width=0.88\linewidth]{figures/report_figures/gd_eta_sweep_vs_muon_best_d128.png}
    \caption{$d=128$. Constant-GD absolute learning-rate sweep compared with the selected constant Muon curve.}
    \label{fig:gd-eta-sweep-d128}
\end{figure}

\begin{figure}[H]
    \centering
    \img[width=0.88\linewidth]{figures/report_figures/normalized_gd_eta_sweep_vs_muon_best_d64.png}
    \caption{$d=64$. Constant-step normalized-GD sweep compared with the selected constant Muon curve.}
    \label{fig:normalized-gd-sweep-d64}
\end{figure}

\begin{figure}[H]
    \centering
    \img[width=0.88\linewidth]{figures/report_figures/normalized_gd_eta_sweep_vs_muon_best_d128.png}
    \caption{$d=128$. Constant-step normalized-GD sweep compared with the selected constant Muon curve.}
    \label{fig:normalized-gd-sweep-d128}
\end{figure}

\end{document}
